\exercisesection{4}

\begin{enumerate}
\def\labelenumi{\arabic{enumi}.}
\tightlist
\item
  Let A be an integrally closed Noetherian domain, and V a vector space of finite rank over the field of fractions of A. Show that, if \((\mathbf{M}_{\lambda})\) is any family of reflexive lattices of V all containing the same lattice N, the lattice \(M = \bigcap_{\lambda} M_{\lambda}\) is reflexive (consider the dual lattices \(M_{\lambda}^{*}\) ).
\end{enumerate}

* 2. Let A be an integrally closed Noetherian domain and E, F two finitely generated A-modules. Suppose that E is torsion-free, that \(\mathrm{Hom}_{\mathrm{A}}(\mathrm{E},\mathrm{F})\) is reflexive and that \(\mathrm{Ext}_{\mathrm{A}}^{1}(\mathrm{E},\mathrm{F}) = 0\) . Show then that F is reflexive. (Prove first that F is torsion-free; if \(\mathrm{T} = \mathrm{F}^{**}/c_{\mathrm{F}}(\mathrm{F})\) , calculate \(\mathrm{Ass}(\mathrm{Hom}_{\mathrm{A}}(\mathrm{E},\mathrm{T}))\) in two ways using the exact sequence:

\[
0 \rightarrow \operatorname{Hom} (\mathrm{E}, \mathrm{F}) \rightarrow \operatorname{Hom} (\mathrm{E}, \mathrm{F} ^ {* *}) \rightarrow \operatorname{Hom} (\mathrm{E}, \mathrm{T}) \rightarrow 0
\]

and Chapter IV, § 1, no. 4, Proposition 10.) *

\begin{enumerate}
\def\labelenumi{\arabic{enumi}.}
\setcounter{enumi}{2}
\item
  Let A be an integrally closed Noetherian domain, K its field of fractions, M a reflexive lattice of a vector space of finite rank over K and L a free lattice of V containing M. Show that there exists a free lattice \(L_{1}\) of V such that \(M = L \cap L_{1}\) . (Consider the finite set I of prime ideals p of A of height 1 such that \(L_{p} \neq M_{p}\) and the principal ideal domain \(S^{-1}A\) , where \(S = \bigcap_{p \in I} (A - p)\) ; show that there exist a free lattice \(L_{0}\) of V such that \(M_{p} = (L_{0})_{p}\) for all \(p \in I\) and an \(s \in S\) such that \(M \subset s^{-1}L_{0} \subset L_{1}\) .)
\item
  Let k be a field and A the polynomial ring \(k[X, Y]\) .
\end{enumerate}

\begin{enumerate}
\def\labelenumi{(\alph{enumi})}
\item
  Let \((e_1, e_2)\) be the canonical basis of the A-module \(A^2\) and let E be the sub-A-module of \(A^2\) generated by \((X - Y)e_1\) , \(e_1 + Xe_2\) and \(e_1 + Ye_2\) ; let F be the monogenous submodule of E generated by \((X - Y)^2e_1\) . Show that the A-module \(M = E / F\) is not the direct sum of its torsion submodule and a torsion-free module.
\item
  Show that the torsion A-module A/AXY is not a direct sum of monogenous submodules of the form \(A/P_{i}^{n_{i}}\) , where the \(P_{i}\) are prime ideals of height 1.
\end{enumerate}

\begin{enumerate}
\def\labelenumi{\arabic{enumi}.}
\setcounter{enumi}{4}
\tightlist
\item
  Let A be an integrally closed Noetherian domain and M a finitely generated A-module. Show that, if a, b are two ideals of A, the A-module
\end{enumerate}

\[
((a \mathbf {M}) \cap (\mathfrak {b} \mathbf {M})) / (a \cap \mathfrak {b}) \mathbf {M}
\]

is pseudo-zero; give an example where it is not zero.

\begin{enumerate}
\def\labelenumi{\arabic{enumi}.}
\setcounter{enumi}{5}
\tightlist
\item
  Let \(k\) be a field, \(\mathbf{B}\) the polynomial ring \(k[\mathbf{X},\mathbf{Y}]\) and \(\mathbf{A}\) the submodule \(k[\mathbf{X}^2,\mathbf{XY},\mathbf{Y}^2]\) of \(\mathbf{B}\) .
\end{enumerate}

\begin{enumerate}
\def\labelenumi{(\alph{enumi})}
\item
  Show that A is an integrally closed Noetherian domain and that B is a finitely generated A-module.
\item
  Show that the ideal \(p = AX^{2} + AXY\) of A is a prime ideal of height 1 (and hence divisorial) but that the B-module \(p \otimes_{A} B\) is a torsion module \(\neq 0\) and hence not reflexive; the ideal pB of B is not divisorial, the canonical mapping \(p \otimes_{A} B \to pB\) is not injective and p is not a flat A-module.
\end{enumerate}

\begin{enumerate}
\def\labelenumi{\arabic{enumi}.}
\setcounter{enumi}{6}
\tightlist
\item
  Let A be an integrally closed Noetherian domain, E a finitely generated torsion-free A-module and \(\mathbf{E}^*\) its dual.
\end{enumerate}

\begin{enumerate}
\def\labelenumi{(\alph{enumi})}
\item
  Show that the canonical homomorphism \(\mathbf{E}^{*}\otimes_{\mathbf{A}}\mathbf{E}\to \mathrm{End}_{\mathbf{A}}(\mathbf{E})\) is a pseudo-isomorphism.
\item
  Deduce from (a) that for E to be a projective A-module, it is necessary and sufficient that \(E^{*} \otimes_{A} E\) be a reflexive A-module. (Note that, if \(E^{*} \otimes_{A} E\) is reflexive, the canonical homomorphism \(E^{*} \otimes_{A} E \to \operatorname{End}_{A}(E)\) is bijective.)
\end{enumerate}

¶* 8. Let A be an integrally closed Noetherian domain and \(\mathbf{M}_1, \mathbf{M}_2\) two finitely generated A-modules.

\begin{enumerate}
\def\labelenumi{(\alph{enumi})}
\item
  Show that the A-modules \(\operatorname{Tor}_i^{\mathbf{A}}(\mathbf{M}_1, \mathbf{M}_2)\) , \(\operatorname{Ext}_{\mathbf{A}}^{i}(\mathbf{M}_1, \mathbf{M}_2)\) are pseudo-zero for \(i \geqslant 2\) (reduce it to the case where \(\mathbf{A}\) is a principal ideal domain).
\item
  If \(\mathbf{M}_1\) is torsion-free, show that \(\operatorname{Tor}_1^{\mathrm{A}}(\mathbf{M}_1, \mathbf{M}_2)\) and \(\operatorname{Ext}_A^1(\mathbf{M}_1, \mathbf{M}_2)\) are pseudo-zero (same method).
\item
  If \(\mathbf{M}_1\) is a torsion A-module, show that (in the notation of no. 5):
\end{enumerate}

\[
\chi (\mathbf {M} _ {1} \otimes_ {\mathbf {A}} \mathbf {M} _ {2}) - \chi (\operatorname{Tor} _ {1} ^ {\mathbf {A}} (\mathbf {M} _ {1}, \mathbf {M} _ {2})) = r (\mathbf {M} _ {2}) \chi (\mathbf {M} _ {1})
\]

\[
\chi (\operatorname{Hom} _ {\mathbf {A}} (\mathbf {M} _ {1}, \mathbf {M} _ {2})) - \chi (\operatorname{Ext} _ {\mathbf {A}} ^ {1} (\mathbf {M} _ {1}, \mathbf {M} _ {2})) = - r (\mathbf {M} _ {2}) \chi (\mathbf {M} _ {1})
\]

\[
\chi (\mathrm{Hom} _ {\mathrm{A}} (\mathrm{M} _ {2}, \mathrm{M} _ {1})) - \chi (\mathrm{Ext} _ {\mathrm{A}} ^ {1} (\mathrm{M} _ {2}, \mathrm{M} _ {1})) = r (\mathrm{M} _ {2}) \chi (\mathrm{M} _ {1}).
\]

\begin{enumerate}
\def\labelenumi{(\alph{enumi})}
\setcounter{enumi}{3}
\tightlist
\item
  Show that for every ordered pair of finitely generated A-modules \(M_{1}\) , \(M_{2}\) :
\end{enumerate}

\[
\begin{array}{r} c (\mathbf {M _ {1}} \otimes_ {\mathbf {A}} \mathbf {M _ {2}}) - c (\mathrm{Tor} _ {\mathbf {1}} ^ {\mathbf {A}} (\mathbf {M _ {1}}, \mathbf {M _ {2}})) = r (\mathbf {M _ {1}}) c (\mathbf {M _ {2}}) + r (\mathbf {M _ {2}}) c (\mathbf {M _ {1}}) \\ c (\mathrm{Hom} _ {\mathbf {A}} (\mathbf {M _ {1}}, \mathbf {M _ {2}})) - c (\mathrm{Ext} _ {\mathbf {A}} ^ {1} (\mathbf {M _ {1}}, \mathbf {M _ {2}})) = r (\mathbf {M _ {1}}) c (\mathbf {M _ {2}}) - r (\mathbf {M _ {2}}) c (\mathbf {M _ {1}}). \end{array}
\]

\begin{enumerate}
\def\labelenumi{\arabic{enumi}.}
\setcounter{enumi}{8}
\tightlist
\item
  Let k be a field and A the polynomial ring \(k[X,Y]\) . On the A-module \(M = A^{2}\) consider the linear form f such that \(f(e_{1}) = X\) , \(f(e_{2}) = Y((e_{1}, e_{2})\) being the canonical basis). Show that the kernel L of f is a monogenous free A-module, but that the quotient M/L (which is isomorphic to an ideal of A) is not reflexive.
\end{enumerate}

¶ 10. Let A be a commutative ring. For every submodule R of a finitely generated free A-module \(L = A^n\) , let \(c_1(R)\) denote the ideal generated by the \(\langle x, x^* \rangle\) , where \(x\) runs through R and \(x^*\) runs through the dual \(L^*\) ; write \(c_k(R) = c_1\left(\operatorname{Im}\left(\bigwedge^k R\right)\right)\) , \(\operatorname{Im}\left(\bigwedge^k R\right)\) being the canonical image of the \(k\) -th exterior power of R in \(\bigwedge^k L\) . If \(R_1 \subset R_2\) are two submodules of L, then \(c_k(R_1) \subset c_k(R_2)\) for all \(k\) .

\begin{enumerate}
\def\labelenumi{(\alph{enumi})}
\item
  Let M be a finitely generated A-module; M is isomorphic to a quotient module L/R, where \(L = A^{n}\) for a suitable n.~Show that the sequence of ideals \((\mathfrak{a}_{k})_{k \geqslant 0}\) such that \(\mathfrak{a}_{k} = \mathfrak{c}_{n-k}(R)\) for \(k \leqslant n\) and \(\mathfrak{a}_{k} = A\) for k \textgreater{} n is independent of the expression of M in the form L/R. (Consider first the case where L/R and L/R' are isomorphic; then note that \(A^{n}/R\) is isomorphic to \(A^{n+h}/(R \times A^{h})\) for all h \textgreater{} 0.) We write \(\mathfrak{d}_{k}(M) = \mathfrak{a}_{k}\) for all \(k \geqslant 0\) and say that these are the determinantal ideals associated with M. Then \(\mathfrak{d}_{k}(M) \subset \mathfrak{d}_{k+1}(M)\) for \(k \geqslant 0\) .
\item
  If \(\mathbf{A}\) is a principal ideal domain and \(\mathbf{M} = \mathbf{L} / \mathbf{R}\) where \(\mathbf{L}\) is free and finitely generated, show that the ideals \(\mathfrak{c}_{k + 1}(\mathbb{R})(\mathfrak{c}_k(\mathbb{R}))^{-1}\) are the invariant factors of \(\mathbf{R}\) in \(\mathbf{L}\) .
\item
  Let \(r = \gamma(\mathbf{M})\) be the least of the cardinals of the systems of generators of \(\mathbf{M}\) and let \(r_0\) be the least integer \(h\) such that \(\mathfrak{d}_h(\mathbf{M}) = \mathbf{A}\) . Show that \(r_0 \leqslant r\) and give an example where \(r_0 < r\) (take \(\mathbf{A}\) to be a Dedekind domain).
\item
  If \(a\) is the annihilator of \(\mathbf{M}\) , show that \(\mathfrak{d}_0(\mathbf{M}) \subset \mathfrak{a}\) and \(\mathfrak{a}^{r - k} \subset \mathfrak{d}_k(\mathbf{M})\) for \(k \leqslant r = \gamma (\mathbf{M})\) .
\item
  Suppose that \(\mathbf{M}\) is a direct sum of submodules \(\mathbf{M}_i\) ( \(1 \leqslant i \leqslant h\) ). Show that
\end{enumerate}

\[
\mathfrak {d} _ {k} (\mathbf {M}) = \sum \mathfrak {d} _ {k _ {1}} (\mathbf {M} _ {1}) \dots \mathfrak {d} _ {k _ {n}} (\mathbf {M} _ {n})
\]

where the sum is over the finite sequences \((k_{i})_{1\leqslant i\leqslant h}\) such that \(\sum_{i = 1}^{h}k_{i} = k.\)

\begin{enumerate}
\def\labelenumi{(\alph{enumi})}
\setcounter{enumi}{5}
\tightlist
\item
  If N is a finitely generated submodule of M, then \(\mathfrak{d}_{k}(\mathbf{M}/\mathbf{N}) \subset \mathfrak{d}_{k}(\mathbf{M})\) for all \(k \geqslant 0\) . Show that:
\end{enumerate}

\[
\sum_ {j = 0} ^ {k} \mathfrak {d} _ {j} (\mathrm{N}) \mathfrak {d} _ {k - j} (\mathrm{M} / \mathrm{N}) \subset \mathfrak {d} _ {k} (\mathrm{M}).
\]

\begin{enumerate}
\def\labelenumi{(\alph{enumi})}
\setcounter{enumi}{6}
\tightlist
\item
  Let K be a field, A the polynomial ring K{[}X, Y, Z{]}, m the maximal ideal AX + AY + AZ and q the ideal generated by X, Y², YZ and Z². Consider the A-module M = A/q and its submodule N = m/q. Show that \(\mathfrak{d}_{0}(M) = \mathfrak{q}, \mathfrak{d}_{1}(M) = A, \mathfrak{d}_{0}(N) = m^{2}\) and \(\mathfrak{d}_{1}(N) = m\) .
\end{enumerate}

\begin{enumerate}
\def\labelenumi{\arabic{enumi}.}
\setcounter{enumi}{10}
\tightlist
\item
  Let A be a Krull domain and M, N two finitely generated lattices in a vector space V of finite rank n over the field of fractions of A. For all \(c \neq 0\) in A such that \(cN \subset M\) , consider the determinantal ideals \(\mathfrak{d}_{k}(M/cN)\) (Exercise 10) and write:
\end{enumerate}

\[
d _ {k} (\mathbf {N}, \mathbf {M}) = \operatorname{div} \left(\mathfrak {d} _ {k} (\mathbf {M} / c \mathbf {N})\right) - (n - k) \operatorname{div} (c).
\]

\begin{enumerate}
\def\labelenumi{(\alph{enumi})}
\item
  Show that \(d_{k}(\mathbf{N}, \mathbf{M})\) does not depend on the choice of c such that \(cN \subset M\) ; \(d_{k}(\mathbf{N}, \mathbf{M})\) is called the determinantal divisor of index k of N with respect to M.
\item
  \(d_{k}(\mathbf{N},\mathbf{M})\geqslant d_{k + 1}(\mathbf{N},\mathbf{M})\) and \(d_{n}(\mathbf{N},\mathbf{M}) = 0\) . Show that, if we write
\end{enumerate}

\[
e _ {k} (\mathrm{N}, \mathrm{M}) = d _ {n - k} (\mathrm{N}, \mathrm{M}) - d _ {n - k + 1} (\mathrm{N}, \mathrm{M}),
\]

then

\[
e _ {k} (\mathbf {N}, \mathbf {M}) \leqslant e _ {k + 1} (\mathbf {N}, \mathbf {M})
\]

for \(1 \leqslant k \leqslant n\) ; the divisors \(e_k(\mathbf{N}, \mathbf{M})\) are called the invariant factors of \(\mathbf{N}\) with respect to \(\mathbf{M}\) (reduce it to the case where \(\mathbf{A}\) is a principal ideal domain).

\begin{enumerate}
\def\labelenumi{(\alph{enumi})}
\setcounter{enumi}{2}
\item
  Show that \(d_0(\mathbf{N}, \mathbf{M}) = \chi(\mathbf{M}, \mathbf{N})\) (no. 5).
\item
  If M is a lattice in V but N is a lattice in a subspace W of V of rank q \textless{} n, the invariant factors of N with respect to \(M \cap W\) are called the invariant factors of N with respect to M. Show how Exercises 8 to 10 and 14 to 16 of Algebra, Chapter VII, §4 may be extended (assuming if need be in certain cases that A is an integrally closed Noetherian domain).
\end{enumerate}

\begin{enumerate}
\def\labelenumi{\arabic{enumi}.}
\setcounter{enumi}{11}
\tightlist
\item
  Let A be an integrally closed Noetherian domain and M, N two torsion-free finitely generated A-modules of the same rank r such that \(N \subset M\) ; let \(j: N \to M\) be the canonical injection.
\end{enumerate}

\begin{enumerate}
\def\labelenumi{(\alph{enumi})}
\tightlist
\item
  Show that, for all \(k\) , \(\bigwedge^{k}j\colon\bigwedge^{k}\mathrm{N}\to\bigwedge^{k}\mathrm{M}\) is pseudo-injective, that \(\operatorname{Coker}\left(\bigwedge^{k}j\right)\) is a torsion A-module and that:
\end{enumerate}

\[
\chi \left(\operatorname{Coker} \left(\bigwedge^ {k} j\right)\right) = \binom {r - 1} {k - 1} \chi (\operatorname{Coker} j).
\]

\begin{enumerate}
\def\labelenumi{(\alph{enumi})}
\setcounter{enumi}{1}
\tightlist
\item
  Using (a), show that, if M is a torsion-free finitely generated A-module, the torsion submodule of \(\bigwedge^{k}\mathbf{M}\) is pseudo-zero for all \(k\) and
\end{enumerate}

\[
c \left(\bigwedge^ {k} \mathbf {M}\right) = \binom {r - 1} {k - 1} c (\mathbf {M}),
\]

where \(\mathbf{M}\) is of rank \(r\) ; in particular, \(c\left(\bigwedge^{r}\mathbf{M}\right) = c(\mathbf{M})\) .

\begin{enumerate}
\def\labelenumi{\arabic{enumi}.}
\setcounter{enumi}{12}
\tightlist
\item
  Let \(k\) be a field and \(A\) the polynomial ring \(k[X, Y]\) .
\end{enumerate}

\begin{enumerate}
\def\labelenumi{(\alph{enumi})}
\item
  Let \(m\) be the maximal ideal \(AX + AY\) of \(A\) ; show that there exists a pseudo-isomorphism of \(m\) to \(A\) , but that there exists no pseudo-isomorphism of \(A\) to \(m\) .
\item
  Let \(\mathfrak{m}'\) be the maximal ideal \(\mathbf{A}(\mathbf{X} - 1) + \mathbf{A}\mathbf{Y}\) of \(\mathbf{A}\) ; show that \(c(\mathfrak{m}) = c(\mathfrak{m}') = 0\) , but that there exists no pseudo-isomorphism of \(\mathfrak{m}\) to \(\mathfrak{m}'\) nor of \(\mathfrak{m}'\) to \(\mathfrak{m}\) .
\item
  Let \(\mathfrak{p} = \mathbf{AX}\) , \(\mathfrak{q} = \mathbf{AY}\) , which are prime ideals of height 1. In the A-module \(\mathbf{L} = \mathbf{A}^2\) , consider the submodules \(\mathbf{M} = \mathfrak{pe}_1 \oplus \mathfrak{qe}_2\) , \(\mathbf{N} = \mathbf{Ae}_1 \oplus \mathfrak{pqe}_2\) ( \(e_1, e_2\) being the vectors of the canonical basis of \(\mathbf{L}\) ). Show that \(\mathbf{M}\) and \(\mathbf{N}\) are isomorphic and that there exist a pseudo-isomorphism from \(\mathbf{L} / \mathbf{M}\) to \(\mathbf{L} / \mathbf{N}\) and a pseudo-isomorphism from \(\mathbf{L} / \mathbf{N}\) to \(\mathbf{L} / \mathbf{M}\) but that there exists no pseudo-isomorphism from \(\mathbf{L}\) to itself mapping \(\mathbf{M}\) to \(\mathbf{N}\) or \(\mathbf{N}\) to \(\mathbf{M}\) (observe that a pseudo-isomorphism from \(\mathbf{L}\) to itself is necessarily an automorphism of \(\mathbf{L}\) , with the aid of Proposition 10).
\end{enumerate}

¶ 14. Let A be a Prüferian domain (§ 2, Exercise 12) and M, N two finitely generated lattices in a vector space V of finite rank n over the field of fractions of A. For all \(c \neq 0\) in A such that \(cN \subset M\) , consider the determinantal ideals \(\mathfrak{d}_{k}(M/cN)\) (Exercise 10) and write:

\[
\mathfrak {d} _ {k} (\mathbf {N}, \mathbf {M}) = c ^ {k - n} \mathfrak {d} _ {k} (\mathbf {M} / c \mathbf {N}).
\]

\begin{enumerate}
\def\labelenumi{(\alph{enumi})}
\item
  Show that \(\mathfrak{d}_k(\mathbf{N},\mathbf{M})\) does not depend on the choice of \(c\) such that \(c\mathbf{N}\subset \mathbf{M}\) ; these are called the determinantal ideals of \(\mathbf{N}\) with respect to \(\mathbf{M}\) .
\item
  \(\mathfrak{d}_k(\mathbf{N},\mathbf{M})\subset \mathfrak{d}_{k + 1}(\mathbf{N},\mathbf{M})\) and \(\mathfrak{d}_n(\mathbf{N},\mathbf{M}) = \mathbf{A}\) . Show that, if we write \(\mathfrak{e}_k(\mathbf{N},\mathbf{M}) = \mathfrak{d}_{n - k}(\mathbf{N},\mathbf{M})(\mathfrak{d}_{n - k + 1}(\mathbf{N},\mathbf{M}))^{-1}\) , the integral ideals \(\mathfrak{e}_k(\mathbf{N},\mathbf{M})\) are such that \(\mathfrak{e}_k(\mathbf{N},\mathbf{M})\supset \mathfrak{e}_{k + 1}(\mathbf{N},\mathbf{M})\) for \(1\leqslant k\leqslant n\) ; the finitely generated ideals \(\mathfrak{e}_k(\mathbf{N},\mathbf{M})\) are called the invariant factors of \(\mathbf{N}\) with respect to \(\mathbf{M}\) (consider the \(\mathrm{A_m}\) -modules \(\mathrm{M_m}\) and \(\mathrm{N_m}\) for every maximal ideal \(\mathfrak{m}\) of \(\mathrm{A}\) ).
\item
  If M is a lattice of V but N is a lattice in a subspace W of V of rank q \textless{} n, the invariant factors of N with respect to M ∩ W are called the invariant factors of N with respect to M. Show how Exercises 8 to 10 and 14 to 16 of Algebra, Chapter VII, § 4 may be extended.
\end{enumerate}

\begin{enumerate}
\def\labelenumi{\arabic{enumi}.}
\setcounter{enumi}{14}
\tightlist
\item
  Let A be the polynomial ring Z{[}X, Y, T, U, V, W{]}, L the free A-module \(A^{4}\) , \((e_{i})_{1 \leq i \leq 4}\) its canonical basis and M the submodule of L generated by the four vectors \(Xe_{1}\) , \(Ye_{2}\) , \(Te_{3} + Ue_{4}\) , \(Ve_{3} + We_{4}\) . Show that
\end{enumerate}

\[
\mathfrak {d} _ {1} (\mathrm{L} / \mathrm{M}) \mathfrak {d} _ {3} (\mathrm{L} / \mathrm{M}) \notin (\mathfrak {d} _ {2} (\mathrm{L} / \mathrm{M})) ^ {2}
\]

(compare with Exercise 14 (b)).

¶ 16. (a) Let A be a Prüferian domain and M a finitely generated lattice in a vector space V of finite rank n over the field of fractions K of A. Show that there exists a basis \((e_{i})_{1 \leqslant i \leqslant n}\) of V and n fractional ideals \(a_{i}\) ( \(1 \leqslant i \leqslant n\) ) such that M is equal to the direct sum of the \(a_{i}e_{i}\) . (Argue by induction on n: if \((u_{i})_{1\leq i\leq n}\) is a basis of V, consider the finitely generated fractional ideal \(b_{1}\) generated by the \(u_{1}\) -coordinates of the elements of M; by considering the A-module \(b_{1}^{-1}M\) , reduce it the case where \(b_{1}=A\) .)

Deduce that, if W is a subspace of V, M ∩ W is a direct factor of M.

\begin{enumerate}
\def\labelenumi{(\alph{enumi})}
\setcounter{enumi}{1}
\tightlist
\item
  Suppose that M is a submodule of \(A^{n}\) equal to the direct sum of the \(a_{i}e_{i}\) , where the \(a_{i}\) are finitely generated (integral) ideals of A and \((e_{i})_{1\leqslant i\leqslant n}\) the canonical basis of \(A^{n}\) . Show that there exists a basis \((u_{i})_{1\leqslant i\leqslant n}\) of \(A^{n}\) and ideals \(b_{i}\) of A such that M is the direct sum of the \(b_{i}u_{i}\) and \(b_{i}\) divides \(b_{i+1}\) for \(1\leqslant i\leqslant n-1\) . (Reduce it to the case where n=2 and \(a_{1}+a_{2}=A\) .)
\end{enumerate}

\begin{enumerate}
\def\labelenumi{\arabic{enumi}.}
\setcounter{enumi}{16}
\tightlist
\item
  Let k be a field, A the polynomial ring \(k[X, Y]\) , L the free A-module \(A^{2}\) and \((e_{1}, e_{2})\) the canonical basis of L. Let M be the submodule of L generated by the two vectors \((X + 1)e_{1} + Ye_{2}, Ye_{1} + Xe_{2}\) . Show that there exists no pseudo-isomorphism from L to itself, whose restriction to M is a pseudo-isomorphism from M to the submodule N of the form \(au + bv\) where \((u, v)\) is a basis of the A-module L and a, b two ideals of A, or whose restriction to N is a pseudo-isomorphism from N to M. (By considering the determinantal ideals (Exercise 10) and noting that M is reflexive, attention may be confined to the case where N would also be reflexive and then necessarily a = A, b = AP, where
\end{enumerate}

\[
\mathrm{P} (\mathrm{X}, \mathrm{Y}) = \mathrm{X} (\mathrm{X} + 1) - \mathrm{Y} ^ {2}
\]

is an extremal element of A; show finally that, for every basis \((u, v)\) of L, M ∩ Au can neither contain Au nor be contained in APu.)

¶ 18. Let A be a Dedekind domain, K its field of fractions and M, N two lattices in a vector space V of finite rank n over K. Let \(\mathfrak{e}_{k} = \mathfrak{e}_{k}(N, M)\) be the invariant factors of N with respect to M (Exercise 14). Show that there exists a basis \((u_{i})_{1 \leqslant i \leqslant n}\) of V such that M is equal to a direct sum \(\bigoplus_{i} a_{i} u_{i}\) , where the \(a_{i}\) are fractional ideals, and N is equal to the direct sum \(\bigoplus_{i} e_{i} a_{i} u_{i}\) . (Use the theory of finitely generated modules over a principal ideal domain (Algebra, Chapter VII, § 4, no. 2, Theorem 1) and the approximation theorem in the unimodular group \(\mathbf{SL}(n, \mathbf{A})\) (§ 2, no. 4).)

Extend this result to the case where A is the union of a right directed family of subrings which are Dedekind domains (for example the integral closure of a Dedekind domain in the algebraic closure of its field of fractions).

\begin{enumerate}
\def\labelenumi{\arabic{enumi}.}
\setcounter{enumi}{18}
\item
  Let A be a Dedekind domain and a, b fractional ideals of A. Show that the A-modules A ⊕ ab and a ⊕ b are isomorphic.
\item
  Let A be a Dedekind domain, P a projective A-module and N a submodule of P. Show that N is projective and a direct sum of modules isomorphic to ideals of A. (Attention may be confined to the case where P is free. Then proceed as in Exercise 16, using a transfinite induction argument based on that of Algebra, Chapter VII, § 3, Theorem 1.)
\item
  Let P be a projective module over a Dedekind domain A. Show that, if P is not finitely generated, it is free. (By virtue of Exercise 20, we may write P as an infinite direct sum \(\bigoplus_{\lambda} (\mathfrak{b}_{\lambda} \oplus \mathfrak{c}_{\lambda})\) , where, for each \(\lambda, \mathfrak{b}_{\lambda}\) and \(\mathfrak{c}_{\lambda}\) are ideals of A. Using Exercise 19, \(\mathbf{P} = \mathbf{L} \oplus \mathbf{Q}\) , where \(\mathbf{L}\) is isomorphic to \(\mathbf{A}^{(I)}\) (I infinite) and \(\mathbf{Q} = \bigoplus_{\alpha \in I} \mathfrak{a}_{\alpha}\) , where the \(\mathfrak{a}_{\alpha}\) are ideals of A. Then apply Exercise 3 of Algebra, Chapter II, § 2.)
\end{enumerate}

¶ 22. Let A be a local ring, m its maximal ideal and E a finitely generated A-module.

\begin{enumerate}
\def\labelenumi{(\alph{enumi})}
\item
  If \(r = \gamma(\mathrm{E})\) is the least of the cardinals of a system of generators of \(\mathbf{E}\) , \(r\) is also the least integer \(h\) such that the determinantal ideal \(\mathfrak{d}_h(\mathbf{E})\) is equal to \(\mathbf{A}\) and also the least of the integers \(k\) such that \(\bigwedge^{k+1} \mathbf{E} = 0\) (note that \(r\) is the rank of \(\mathbf{E}/\mathfrak{m}\mathbf{E}\) over \(\mathbf{A}/\mathfrak{m}\) ).
\item
  If we write \(e(E) = \mathfrak{d}_{r-1}(E)\) , show that \(\bigwedge^{r} E\) is isomorphic to \(A / e(E)\) . For an ideal \(a\) of \(A\) to contain \(e(E)\) , it is necessary and sufficient that \(E / aE\) be a free \((A / a)\) -module (reduce it to the case where \(a = 0\) ).
\item
  If \(\mathfrak{e}(\mathbf{E})\) is a principal ideal \(A\alpha\) , show that E contains a direct factor isomorphic to \(A/A\alpha\) (write E as a quotient of \(A^{r}\) by a submodule R and note that there is in R an element of the form \(\alpha y\) , where y is an element of a basis of \(A^{r}\) ).
\item
  Let \(\mathfrak{b}_h^{\prime}(\mathbf{E})\) be the annihilator of \(\bigwedge^{h+1}\mathbf{E}\) for \(0 \leqslant h \leqslant r - 1\) . Show that the following conditions are equivalent:
\end{enumerate}

\((\alpha)\) The ideals \(\mathfrak{d}_h(\mathbf{E})\) ( \(0 \leqslant h \leqslant r - 1\) ) are principal.

(β) The ideals \(\mathfrak{d}_h^\prime (\mathrm{E})\) \((0\leqslant h\leqslant r - 1)\) are principal.

\((\gamma)\) E is the direct sum of r modules isomorphic to A/A \(\lambda_{h}\) , where \(\lambda_{h+1}\) divides \(\lambda_{h}\) for \(0 \leqslant h \leqslant r - 1\) .

Moreover, if these conditions hold, then \(\mathfrak{d}_h'(\mathrm{E}) = \mathrm{A}\lambda_h\) for \(0 \leqslant h \leqslant r - 1\) and \(\mathfrak{d}_h(\mathrm{E}) = \lambda_h\lambda_{h+1} \ldots \lambda_{r-1}\mathrm{A}\) for \(0 \leqslant h \leqslant r - 1\) . (Proceed by induction on \(r\) , using (c).)

\begin{enumerate}
\def\labelenumi{(\alph{enumi})}
\setcounter{enumi}{4}
\tightlist
\item
  Suppose further that A is an integral domain. Show that, if p is a prime ideal of A such that E/pE is a free (A/p)-module and \(E_{p}\) a free \(A_{p}\) -module, then E is a free A-module (using (b), show that \(\gamma(E_{p}) = \gamma(E)\) and that \(e(E_{p}) = e(E)\) ).
\end{enumerate}

\begin{enumerate}
\def\labelenumi{\arabic{enumi}.}
\setcounter{enumi}{22}
\item
  Let A be a ring, E a finitely generated A-module, r the least of the integers k such that \(\bigwedge^{k+1}E=0\) and \(e(E)\) the annihilator of \(\bigwedge^{r}E\) . Show that, for every ideal \(a\supset e(E)\) in A, the \((A/a)\) -module E/aE is flat (reduce it to the case where A is a local ring and use Exercise 22 (b)). Deduce that, if r is the Jacobson radical of A and if, for every maximal ideal m of A, the rank of the vector \((A/m)\) -space E/mE is the same, then E/rE is a flat \((A/r)\) -module.
\item
  Let A be an integrally closed Noetherian domain. For a lattice M (with respect to A) to be reflexive, it is necessary and sufficient that it satisfy the following condition: for every ordered pair \((a, b)\) of elements of A such that \(a \neq 0\) and the homothety of ratio b on A/aA is injective, the homothety of ratio b on M/aM is then injective. (To see that the condition is necessary, consider two elements x, y of M such that \(ax + by = 0\) and show that, for all \(p \in P(A)\) , \(y \in aM_{p}\) . To show that the condition is sufficient, use criterion (c) of Theorem 2 and observe (using Proposition 8 of §1, no. 4) that the hypothesis may be written as
\end{enumerate}

\[
\operatorname{Ass} (\mathbf {M} / a \mathbf {M}) = \operatorname{Ass} (\mathbf {A} / a \mathbf {A})
\]

for all \(a \neq 0\) in \(\mathbf{A}\) and that for every module \(\mathbf{E}\) such that \(\mathbf{M} \subset \mathbf{E} \subset \mathbf{V}\) there exists \(a \neq 0\) in \(\mathbf{A}\) such that \(a\mathbf{M} \subset a\mathbf{E} \subset \mathbf{M}\) .

\begin{enumerate}
\def\labelenumi{\arabic{enumi}.}
\setcounter{enumi}{24}
\tightlist
\item
  Let \(A = \bigoplus_{n \geq 0} A_n\) be a Noetherian graded ring with positive degrees. Show that, if A is a factorial domain, then \(A_0\) is a factorial domain and the \(A_n\) are reflexive. (To show that \(A_0\) is a factorial domain, use criterion (c) of §3, no. 2, Theorem 1; to show that the \(A_n\) are reflexive, use Exercise 24.)
\end{enumerate}

¶ 26. Let A be a Noetherian ring and M a finitely generated A-module. For the symmetric algebra S(M) to be a factorial domain, it is necessary and sufficient that A be a factorial domain and that the \(S^{n}(M)\) be reflexive A-modules. (To see that the condition is sufficient, observe first that, if T = A - \{0\}, S(M) is identified with a subring of \(T^{-1}S(M)\) ; then show that every extremal element p of A is extremal in S(M), by reducing it to proving that, if p divides a product xy of two homogeneous elements in S(M), it divides one of them; finally, apply Proposition 3 of § 3, no. 4.)

\specialchapter{Historical note (Chapters I to VII)}

(Numbers in brackets refer to the bibliography at the end of this Note.)

``Abstract'' commutative algebra is a recent creation but its development can only be understood as a function of that of the theory of algebraic numbers and algebraic geometry, which gave birth to it.

It has been conjectured without too much improbability that the famous ``proof'' Fermat claimed to possess of the impossibility of the equation \(x^{p} + y^{p} = z^{p}\) for p an odd prime and x, y, z integers \(\neq 0\) depended on the decomposition

\[
(x + y) (x + \zeta y) \dots (x + \zeta^ {p - 1} y) = z ^ {p}
\]

in the ring \(Z[\zeta]\) (where \(\zeta \neq 1\) is a \(p\) -th root of unity) and on a divisibility argument in this ring, assuming it to be a principal ideal domain. In any case an analogous argument is found outlined by Lagrange ({[}2{]}, vol.~II, p.~531); it is by arguments of this type, with certain variations (notably changes of variable aimed at lowering the degree of the equation) that Euler ({[}1{]}, vol.~I, p.~488) (*) and Gauss ({[}3{]}, vol.~II, p.~387) show Fermat's Theorem for \(p = 3\) , Gauss (loc. cit.) and Dirichlet ({[}4{]}, vol.~I, p.~42) for \(p = 5\) and Dirichlet the impossibility of the equation \(x^{14} + y^{14} = z^{14}\) ({[}4{]}, vol.~I, p.~190). Finally, in his first research on the theory of numbers, Kummer believed he had obtained in this way a general proof and it was no doubt this mistake (which Dirichlet pointed out to him) that led him to his study of the arithmetic of cyclotomic fields from which he was eventually to succeed in deducing a correct version of his proof for prime numbers \(p < 100\) {[}7d{]}.

On the other hand, the celebrated memoir of Gauss of 1831 on biquadratic residues, whose results are derived from a detailed study of divisibility in the ring \(\mathbf{Z}[i]\) of ``Gaussian integers'' ({[}3{]}, vol.~II, p.~109), showed clearly the interest that the extension of divisibility to algebraic numbers could hold out for the classical problems of the theory of numbers (*); therefore it is not surprising that between 1830 and 1850 this theory was the subject of numerous works by German mathematicians, first Jacobi, Dirichlet and Eisenstein, and then, a little later, Kummer and his pupil and friend Kronecker. We shall not speak here of the theory of units, which is too specialized a branch of the theory of numbers, where progress was very rapid, Eisenstein obtaining the structure of the group of units for cubic fields and Kronecker for cyclotomic fields, just before Dirichlet in 1846 ({[}4{]}, vol.~I, p.~640) proved the general theorem, at which Hermite had almost arrived independently ({[}8{]}, vol.~I, p.~159). The question (central to the whole theory) of decomposition into prime factors appeared much more difficult. Since Lagrange had given examples of numbers of the form \(x^{2} + \mathrm{D}y^{2}(x,y,\mathrm{D}\) integers) with divisors which are not of the form \(m^2 +\mathrm{D}n^2\) ({[}2{]}, vol.~II, p.~465), it was effectively known that the ring \(\mathbf{Z}[\sqrt{-\mathrm{D}} ]\) could not in general be expected to be a principal ideal domain and Euler's temerity was followed by considerable circumspection; when Dirichlet, for example, proves that the relation \(p^2 -5q^2 = r^5\) ( \(p,q,r\) integers) is equivalent to

\[
p + q \sqrt {5} = (x + y \sqrt {5}) ^ {5}
\]

for integers x, y, he restricts himself to pointing out that ``there are analogous theorems for many other prime numbers {[}than 5{]}'' ({[}4{]}, vol.~I, p.~31). In Gauss's memoir of 1831 and the work of Eisenstein on cubic residues {[}6a{]}, it is certainly true that there are advanced studies on arithmetic in the principal ideal domains Z{[}i{]} and Z{[}ρ{]} (ρ = (-1 + i√3)/2, a cube root of unity) in perfect analogy with the theory of rational integers and in these examples at least the close connection between arithmetic in quadratic fields and the theory of binary quadratic forms developed by Gauss was very apparent; but the general case lacked a ``dictionary'' which would have allowed quadratic fields to be treated by a simple translation from Gauss's theory (†).

In fact, it is not for quadratic fields but for cyclotomic fields (and for reasons which will only appear clearly much later (cf.~p.~585)) that the problem was first solved. From 1837 onwards, Kummer, originally an analyst, turns to the arithmetic of cyclotomic fields which was to occupy him almost exclusively for 25 years. Like his predecessors, he studies divisibility in the ring Z{[}ζ{]}, where ζ is a p-th root of unity ≠1 (p an odd prime); he quickly sees that here also rings are encountered which are not principal ideal domains, blocking all progress in the extension of the laws of arithmetic {[}7a{]} and it is only in 1845, after 8 years' efforts, that the light dawns, thanks to his definition of ``ideal numbers'' ({[}7c{]} and {[}7d{]}).

What Kummer does amounts exactly, in modern language, to defining the valuations on the field \(\mathbf{Q}(\zeta)\) : they are in one-to-one correspondence with his ``ideal prime numbers'', the ``exponent'' with which such a factor appears in the ``decomposition'' of a number \(x \in \mathbf{Z}[\zeta]\) is just the value at \(x\) of the corresponding valuation. As the conjugates of \(x\) also belong to \(\mathbf{Z}[\zeta]\) and their product \(\mathrm{N}(x)\) (the ``norm'' of \(x(*)\) ) is a rational integer, the ``ideal prime factors'' to be defined must also be ``factors'' of the rational prime numbers and in order to define them it was sufficient just to say what were the ``ideal prime divisors'' of a prime number \(q \in \mathbf{Z}\) . For \(q = p\) Kummer had already effectively proved {[}7a{]} that the principal ideal \((1 - \zeta)\) was prime and that its \((p - 1)\) -th power was the principal ideal \((p)\) ; this case therefore raised no new problem. For \(q \neq p\) the idea which seems to have guided Kummer is to replace the cyclotomic equation \(\Phi_p(z) = 0\) by the congruence \(\Phi_p(u) \equiv 0 \pmod{q}\) , in other words to decompose the cyclotomic polynomial \(\Phi_p(\mathbf{X})\) over the field \(\mathbf{F}_q\) and to associate with each irreducible factor of this polynomial an ``ideal prime factor'' A simple case (explicitly mentioned in the Note {[}7b{]} where Kummer announces his results without proof) is that where \(q \equiv 1 \pmod{p}\) ; if \(q = mp + 1\) and \(\gamma \in \mathbf{F}_q\) is a primitive \((q - 1)\) -th root of 1, then, in \(\mathbf{F}_q[\mathbf{X}]\) ,

\[
\Phi_ {p} (\mathbf {X}) = \prod_ {k = 1} ^ {p - 1} \left(\mathbf {X} - \gamma^ {k m}\right)
\]

since \(\gamma^{pm} = 1\) . Then associating with each factor \(\mathbf{X} - \gamma^{km}\) an ``ideal prime factor'' \(\mathfrak{q}_k\) of \(q\) , Kummer says that an element \(x \in \mathbf{Z}[\zeta]\) , of which P is the minimal polynomial over \(\mathbf{Q}\) , is divisible by \(\mathfrak{q}_k\) if in \(\mathbf{F}_q\) \(\mathrm{P}(\gamma^{km}) = 0\) ; to sum up, in modern language, he writes the quotient ring \(\mathbf{Z}[\zeta]/q\mathbf{Z}[\zeta]\) as a direct composition of fields isomorphic to \(\mathbf{F}_q\) . For \(q \not\equiv 1 \pmod{p}\) , the irreducible factors of \(\Phi_p(\mathbf{X})\) in \(\mathbf{F}_q[\mathbf{X}]\) are no longer of first degree and it would therefore be necessary to substitute for X in \(\mathrm{P}(\mathbf{X})\) ``Galois imaginary'' roots of the factors of \(\Phi_p\) in \(\mathbf{F}_q[\mathbf{X}]\) . Kummer avoids this difficulty by passing, as we would say today, to the decomposition field K of \(q\) ; if \(f\) is the least integer such that \(q^f \equiv 1 \pmod{p}\) and \(p - 1 = ef\) , K is just the subfield of \(\mathbf{Q}(\zeta)\) consisting of the invariants of the subgroup of order \(f\) of the Galois group (cyclic of order \(p - 1\) ) of \(\mathbf{Q}(\zeta)\) over \(\mathbf{Q}\) ; in other words it is the unique subfield of \(\mathbf{Q}(\zeta)\) which is of degree \(e\) over \(\mathbf{Q}\) ; it had been well known since Gauss's Disquisitiones, being generated by the ``periods''

\[
\eta_ {k} = \zeta_ {k} + \zeta_ {k + f} + \zeta_ {k + 2 f} + \dots + \zeta_ {k + (e - 1) f}
\]

\((0 \leqslant k \leqslant e - 1, \zeta_{\nu} = \zeta^{g^{\nu}} \text{ where } g \text{ is a primitive root of the congruence } z^{p-1} \equiv 1 (\text{mod. } p))\) , which form a normal basis for it. If \(R(X)\) is the minimal polynomial (monic and with rational integer coefficients) of any of these ``periods'' \(\eta\) , Kummer, starting from Gauss's formulae, proves that, over the field \(F_q\) , \(R(X)\) also decomposes into distinct factors of the first degree \(X - u_j\) ( \(1 \leqslant j \leqslant e\) ) and it is with each of the \(u_j\) that he now associates an ``ideal prime factor'' \(q_j\) . To define ``divisibility by \(q_j\) '', Kummer writes every \(x \in Z[\zeta]\) in the

form \(x = \sum_{k=0}^{\infty} \zeta^k y_k\) , where each \(y_k \in \mathbf{K}\) may itself be written uniquely as a polynomial of degree \(\leqslant e - 1\) in \(\eta\) with rational integer coefficients; he says that \(x\) is divisible by \(q_j\) if and only if, when \(u_j\) is substituted for \(\eta\) in each of the \(y_k\) , the elements of \(\mathbf{F}_q\) obtained are all zero. But it was also necessary to define the ``exponent'' of \(q_j\) in \(x\) . For this, Kummer introduces what we would now call a uniformizer for \(q_j\) , that is an element \(\rho_j \in \mathbf{K}\) such that \(N(\rho_j) \equiv 0 \pmod{q}\) , \(N(\rho_j) \not\equiv 0 \pmod{q^2}\) and finally such that \(\rho_j\) is divisible by \(q_j\) (in the sense defined above) but by none other of the ideal factors \(\neq q_j\) of \(q\) . The existence of such a \(\rho_j\) had effectively been proved by Kronecker in his dissertation the previous year ({[}9a{]}, p.~23); then writing \(\rho_j' = N(\rho_j)/\rho_j\) , Kummer says that the exponent of \(q_j\) in \(x\) is equal to \(h\) if \(x\rho_j'^h \equiv 0 \pmod{q^h}\) but \(x\rho_j'^{h+1} \not\equiv 0 \pmod{q^{h+1}}\) ; he begins of course by proving that the relation \(x\rho_j' \equiv 0 \pmod{q}\) is equivalent to the fact that \(x\) is divisible by \(q_j\) (in the above sense). Once these definitions were made, the extension to \(Z[\zeta]\) of the usual laws of divisibility for ``ideal numbers'' no longer offered serious difficulty; and from his first memoir {[}7c{]} Kummer could even, using Dirichlet's ``box method'', show that the ``classes'' and ``ideal factors'' were finite in number (*).

We shall not pursue the history of Kummer's later works on cyclotomic fields, concerning the determination of the class number and the application to the proof of Fermat's theorem in certain cases. Let us just mention the way in which, in 1859, he extends his method to obtain (at least partially) the ``ideal prime numbers'' in a ``Kummerian field'' \(\mathbf{Q}(\zeta, \mu)\) , where \(\mu\) is a root of the irreducible polynomial \(\mathrm{P}(\mathbf{X}) = \mathbf{X}^p - \alpha\) , where \(\alpha \in \mathbf{Z}[\zeta]\) {[}7e{]}. It is interesting that Kummer envisages the problem precisely by considering \(\mathbf{Q}(\zeta, \mu)\) as a cyclic extension of the field \(\mathbf{Q}(\zeta)\) taken as ``base field'' ( \(\dagger\) ): he starts with an ``ideal prime number'' \(q\) of \(\mathbf{Z}[\zeta]\) which he assumes divides neither \(p\) nor \(\alpha\) and this time he examines (in modern terms) the polynomial \(\overline{\mathrm{P}}(\mathbf{X}) = \mathbf{X}^p - \bar{\alpha}\) in the residue field \(k\) of the valuation on \(\mathbf{Q}(\zeta)\) corresponding to \(q\) ( \(\bar{\alpha}\) being the canonical image of \(\alpha\) in \(k\) ). As \(\mathbf{Q}(\zeta)\) is the field of the \(p\) -th roots of unity, \(\overline{\mathrm{P}}\) is, either irreducible over \(k\) , or the product of factors of the first degree. In the first case, Kummer says that \(q\) remains prime in \(\mathbf{Z}[\zeta, \mu]\) ; in the second, he introduces elements \(w_i\) ( \(1 \leqslant i \leqslant p\) ) of \(\mathbf{Z}[\zeta]\) whose images in \(k\) are the roots of \(\overline{\mathrm{P}}\) and he associates with each index \(i\) an ideal prime factor \(\mathfrak{r}_i\) of \(q\) ; then writing \(\mathrm{W}_i(\mathbf{X}) = \prod_{j \neq i} (\mathbf{X} - w_j)\) , he says that, for a polynomial \(f\) with coefficients in \(\mathbf{Z}[\zeta], f(\mu)\) contains the factor \(\mathfrak{r}_i m\) times if

\[
f (w _ {i}) \mathrm{W} _ {i} ^ {m} (w _ {i}) \equiv 0 \quad (\mathrm{mod.} q ^ {m})
\]

but

\[
f (w _ {i}) \mathrm{W} _ {i} ^ {m + 1} (w _ {i}) \not \equiv 0 \pmod {q ^ {m + 1}}.
\]

To sum up, he obtains in this way the valuations on \(\mathbf{Q}(\zeta, \mu)\) which are unramified over \(\mathbf{Q}\) , which are sufficient for the applications he has in view.

\[
\begin{array}{c c c} \ast & \ast & \ast \\ \hline \end{array}
\]

Kummer had had the chance to meet, in studying particular fields to which his research on Fermat's Theorem had led at first, a number of fortuitous circumstances which made their study much more accessible. The extension to the general case of Kummer's results presented considerable difficulties and was to take years of effort.

With Kronecker and Dedekind, who play the principal roles there, the history of the theory of algebraic numbers, during the 40 years following Kummer's discovery, is not dissimilar (but happily without the same acrimony) to that of the rivalry of Newton and Leibniz 180 years earlier concerning the invention of Infinitesimal Calculus. Pupil and later colleague of Kummer in Berlin, Kronecker (whose thesis, as we have seen, had served as an essential point in Kummer's theory) was greatly interested in ``ideal numbers'' with the aim of applying them to his own research; and we admire his astonishing penetration when we see him, as early as 1853 ({[}9b{]}, p.~10) announce the general theorem on the structure of Abelian extensions of \(\mathbf{Q}\) and, what is perhaps still more remarkable, create, in the years which follow, the theory of complex multiplication and discover the first germ of class field theory ({[}9c{]} and {[}9d{]}). A letter from Kronecker to Dirichlet in 1857 ({[}9{]}, vol.~5, pp.~418--421) shows that he already possessed at that time a generalization of Kummer's theory, which moreover Kummer himself confirms in one of his own works ({[}7e{]}, p.~57) and Kronecker will make many an allusion to this theory in his memoirs between 1860 and 1880 (*).

But although at that time none of the mathematicians of the German school of the Theory of Numbers was unaware of the existence of these works of Kronecker, the latter seems only to have communicated the principles of his methods to a restricted circle of friends and pupils and when he finally decided to publish them in his memoir of 1881 on the discriminant {[}9e{]} and above all in his great ``Festschrift'' of 1882 {[}9f{]}, Dedekind could not refrain from expressing his surprise ({[}10{]}, vol.~III, p.~427), having imagined the processes were completely different, from the echoes he had heard ({[}10{]}, vol.~III, p.~287). Kronecker moreover was far from possessing to the same degree Dedekind's remarkable gifts of exposition and clarity and it is therefore not surprising that it is chiefly the methods of the latter, already published in 1871, which have formed the framework of the theory of algebraic numbers; however interesting it may be, Kronecker's method of the ``adjunction of indeterminates'', where the Theory of Numbers is concerned, is scarcely more in our eyes than a variant of Dedekind's (cf.~Chapter VII, § 1, Exercise 31) and it is chiefly in another direction, oriented towards Algebraic Geometry, that Kronecker's ideas acquire all their importance for the history of Commutative Algebra, as we shall see later.

For reasons which could only clearly be seen much later, a first preliminary to any attempt at a general theory was of course the clarification of the notion of algebraic integer. This is obtained about 1845--50, although it is difficult enough to date its appearance precisely; it seems probable that it is the idea of a system stable under addition and multiplication (or, more precisely, what we now call a Z-algebra of finite rank) which, more or less consciously, led to the general definition of algebraic integers: in fact this definition is inevitably hit upon when a Z-algebra of the form Z{[}θ{]} is restricted to being of finite rank, by analogy with the ring Z{[}ζ{]} generated by a root of unity, which was always at the centre of arithmeticians' preoccupations at this time. At any rate, when, independently, Dirichlet ({[}4{]}, vol.~I, p.~640), Hermite ({[}8{]}, vol.~I, pp.~115 and 146) and Eisenstein ({[}6c{]}, p.~236) introduce the notion of algebraic integer, they do not appear to consider that they are dealing with a new concept nor to judge that it will be useful to make a detailed study thereof; only Eisenstein shows effectively (loc. cit.) that the sum and product of two algebraic integers are algebraic integers, without moreover claiming that this result is original.

A much more subtle point was the determination of the rings in which a generalization of Kummer's theory could be expected. The latter, in his first note {[}7b{]}, does not hesitate to affirm that he can regain by his method Gauss's theory of binary quadratic forms by considering the rings \(Z[\sqrt{D}]\) (D an integer); he never developed this idea, but it certainly seems that neither he nor any one else before Dedekind perceived that unique composition into ``ideal'' prime factors is impossible in the ring \(Z[\sqrt{D}]\) when \(D \equiv 1 \pmod{4}\) (although the example of the cube roots of unity showed that the ring \(Z[\rho]\) considered from the time of Gauss is distinct from \(Z[\sqrt{-3}]\) ) (*). Before Dedekind and Kronecker, the only rings studied are always of the type \(Z[\theta]\) or sometimes certain particular rings of the type \(Z[\theta, \theta']\) (†). As far as Kronecker is concerned, it is possible that the idea of considering all the integers of an algebraic extension was first suggested to him by the study of the field of algebraic functions, where this ring arises naturally as the set of functions which are ``finite at infinite distance''; in any case he insists in his memoir of 1881 on the discriminant (written and announced at the Academy of Berlin as early as 1862) on this characterization of the ``integers'' in these fields {[}9e{]}. Dedekind gives no indication as to the origin of his own ideas on this point, but in his very first publication on number fields in 1871 the ring of all integers of such a field plays a capital role in his theory; it is also Dedekind who clarifies the relation between such a ring and its subrings with the same field of fractions, by the introduction of the notion of conductor {[}10c{]}.

But here was not the only difficulty. To generalize the ideas of Kummer, it was necessary first to get rid of passing via the decomposition field, which naturally could have no analogue in the case of a non-Abelian field. This detour seems moreover at first sight very surprising and artificial, for, starting with the irreducible polynomial \(\Phi_p(\mathbf{X})\) of \(\mathbf{Z}[\mathbf{X}]\) , one may wonder why Kummer does not push his ideas to their logical conclusion and what prevents him from using the theory of ``Galois imaginary'' numbers which were well known at the time. The obstacle comes more clearly to light in an unfortunate attempt at generalization made as early as 1865 by Selling, a pupil of Dedekind: given an irreducible polynomial \(\mathbf{P} \in \mathbf{Z}[\mathbf{X}]\) , Selling decomposes the corresponding polynomial \(\overline{\mathbf{P}}(\mathbf{X})\) into irreducible factors in \(\mathbf{F}_q[\mathbf{X}]\) ; the roots of this polynomial therefore belong to a finite extension \(\mathbf{F}_r\) of \(\mathbf{F}_q\) ; but Selling, in order to define in Kummer's way the exponent of an ``ideal prime factor'' of \(q\) in an integer of the splitting field of \(\mathrm{P}(\mathrm{X})\) , does not hesitate to speak, in the field \(\mathbf{F}_r\) , of congruences modulo a power of \(q\) ({[}11{]}, p.~26); and a little later when he tries to approach the question of ramification, he ``adjoins'' to \(\mathbf{F}_r\) ``imaginary roots'' of an equation of the form \(x^h = q\) ({[}11{]}, p.~34). Clearly these bold steps (which would be justified were the finite field \(\mathbf{F}_q\) replaced by the \(q\) -adic field) could at that time only end in nonsense. Fortunately, Dedekind in 1857 {[}10a{]}, under the name of ``theory of higher congruences'', took up again in another form the theory of finite fields (*): he interprets the elements of the latter as ``residues'' of the polynomials of \(\mathbf{Z}[\mathbf{X}]\) with respect to a ``double modulus'' consisting of the linear combinations with coefficients in \(\mathbf{Z}[\mathbf{X}]\) of a prime number \(p\) and an irreducible monic polynomial \(\mathrm{P} \in \mathbf{Z}[\mathrm{X}]\) (which is no doubt for him, as for Kronecker, the origin of the general idea of module at which they were to arrive independently a little later). According to his own testimony ({[}10d{]}, p.~218) it seems that Dedekind had begun by attacking the problem of the ``ideal factors'' of \(p\) in a field \(\mathbf{Q}(\xi)\) , where \(\mathrm{P} \in \mathbf{Z}[\mathrm{X}]\) is the minimal polynomial of \(\xi\) , as follows (certainly at least in the ``unramified'' case, that is when the polynomial \(\overline{\mathbf{P}}\) in \(\mathbf{F}_p[\mathbf{X}]\) corresponding to \(\mathrm{P}\) has no multiple root): he writes, in \(\mathbf{Z}[\mathbf{X}]\) ,

\[
\mathbf {P} = \mathbf {P} _ {1} \mathbf {P} _ {2} \dots \mathbf {P} _ {h} + p. \mathbf {G}
\]

where the \(\overline{\mathbf{P}}_i\) are irreducible and distinct in \(\mathbf{F}_p[\mathbf{X}]\) ; it may be assumed that \(\mathbf{G}\) is not divisible (in \(\mathbf{Z}[\mathbf{X}]\) ) by any of the \(\mathbf{P}_i\) and for all \(i\) he writes \(\mathbf{W}_i = \prod_{j\neq i}\mathbf{P}_j\) ; then, if \(f\in \mathbf{Z}[\mathbf{X}]\) , it will be said that \(f(\xi)\) contains the ``ideal factor'' \(\mathfrak{p}_i\) of \(p\)

corresponding to \(\mathbf{P}_i\) \(k\) times if

\[
f \mathrm{W} _ {i} ^ {k} \equiv 0 \quad (\text { modd. } p ^ {k}, \mathrm{P})
\]

and

\[
f \mathrm{W} _ {i} ^ {k + 1} \not \equiv 0 \quad (\text { modd. } p ^ {k + 1}, \mathrm{P}).
\]

The relationship with the method followed by Kummer for ``Kummerian fields'' is here manifest and Kummer's original definition for cyclotomic fields can, in this way, easily be recovered (see for example the work of Zolotareff {[}14{]} who, at first independently of Dedekind, developed these ideas a little later).

However, neither Dedekind nor Kronecker who appears also to have made analogous attempts, could progress further in this direction, both of them halted by the difficulties presented by ramification ({[}10d{]}, p.~218 and {[}9f{]}, p.~325) (*). If the ring of integers A of the number field K under consideration admits a basis (over Z) consisting of the powers of the same integer \(\theta\) , it is not difficult to generalize the above method for ramified prime numbers in \(\mathbf{Z}[\theta]\) (as Zolotareff indicates (loc. cit.)). But there are fields K where no basis of this type exists in the ring A; and Dedekind even finished by discovering that there are cases where certain prime numbers \(p\) (the ``extraordinary factors of the discriminant'' of the field K) are such that, for all \(\theta \in \mathrm{A}\) , applying the above method of the minimal polynomial of \(\theta\) over Q leads to attributing to \(p\) multiple ideal factors when in fact \(p\) is unramified in A ( \(\dagger\) ); he admits that he was held up for a long time by this unforeseen difficulty, before managing to surmount it by the creation from scratch of the theory of modules and ideals, in a masterly exposition (and already in a wholly modern style, in contrast with the discursive style of his contemporaries) in what is without doubt his masterpiece, the famous ``11th supplement'' to Dirichlet's book on the Theory of Numbers {[}10f{]}. This work saw three successive versions, but already in the first (published as the ``10th supplement'' to the second edition of Dirichlet's book in 1871 {[}4 bis{]}) the essentials of the method are present and almost at one stroke the theory of algebraic numbers passes from sketches and earlier gropings to a fully mature discipline already possessing its essential tools: from the beginning, the ring of all integers of a number field is placed at the centre of the theory; Dedekind proves the existence of a basis of this ring over Z and deduces from it the definition of the discriminant of the field as the square of the determinant formed by the elements of a basis of the ring of integers and their conjugates; however he only gives in the 11th supplement the characterization of ramified prime numbers (as prime factors of the discriminant) for quadratic fields ({[}10f{]}, p.~202), whereas he was in possession of the general theorem from 1871 onwards (*). The central result of the work is the existence and uniqueness theorem for the decomposition of ideals into prime factors, for which Dedekind starts by developing an elementary theory of ``modules''; in fact, in the 11th supplement, he reserves this name for sub-Z-modules of a number field, but the conception he forms of them and the results he shows are already expounded in a way which is immediately applicable to general modules (†); amongst other things must be noted, as early as 1871, the introduction of the notion of ``transporter'' which plays an important role (as well as the ``ascending chain condition'') in the first proof of the unique factorization theorem. In the two following editions, Dedekind was also to give two other proofs of this theorem which he justly considered as the cornerstone of his theory. It should be noted here that it is in the third proof that fractional ideals are made use of (already introduced as early as 1859 by Kummer for cyclotomic fields) and the fact that they form a group is established; we shall return later to the second proof (p.~594).

All these results (except for the language) were no doubt already known to Kronecker about 1860 as particular cases of his more general conceptions of which we speak later (whereas Dedekind recognizes that he only surmounted the last difficulties of his theory in 1869--70 ({[}10e{]}, p.~351)) ( \(\ddagger\) ); as far as number fields are concerned, it must in particular be underlined that, already at this time, Kronecker knew that the whole theory is applicable without essential change starting with a ``base field'' \(k\) which is itself a number field (other than \(\mathbf{Q}\) ), a point of view to which the theory of complex multiplication led naturally; he had thus recognized, for certain fields \(k\) , the existence of algebraic extensions \(\mathbf{K} \neq k\) unramified over \(k\) ({[}9f{]}, p.~269), a fact which cannot hold for \(k = \mathbf{Q}\) (as follows from minorations of Hermite and Minkowski for the discriminant). Dedekind was never to develop this last point of view (although he indicates its possibility in his memoir of 1882 on the different) and the first systematic exposition of ``relative field'' theory is due to Hilbert {[}16d{]}.

Finally, in 1882 {[}10e{]}, Dedekind completes the theory by introducing the different, which gives him a new definition of the discriminant and allows him to define precisely the exponents of the ideal prime factors in the decomposition of the latter. It is also about this time that he becomes interested in the particular features presented by Galois extensions, introducing the notions of decomposition group and inertia group (in his memoir {[}10g{]} which was only published in 1894) and even (in papers not published during his lifetime ({[}10{]}, vol.~II, pp.~410--411)) a sketch of ramification groups, which Hilbert (independently of Dedekind) was to develop a little later ({[}16c{]} and {[}16d{]}).

Thus, about 1895, the theory of algebraic numbers had completed the first stage of its development; the tools forged during this formation period will allow it almost immediately to enter the next stage, general class field theory (or, what amounts to the same, the theory of Abelian extensions of number fields) which carries on to our own day and which we shall not describe here. From the point of view of Commutative Algebra, it may be said that at the same time the history of Dedekind domains is practically completed, setting aside their axiomatic characterization, and also the structure of finitely generated modules over these domains (which, in the case of number fields, will only be substantially elucidated by Steinitz in 1912 {[}20b{]}) (*).

\begin{center}* * *\end{center}

The later progress in Commutative Algebra arises chiefly from quite different problems, issuing from Algebraic Geometry (which will moreover influence the Theory of Numbers directly even before the ``abstract'' developments of the present period).

We shall not concern ourselves here with the detailed history of Algebraic Geometry which, until the death of Riemann, scarcely touches our subject. Let it suffice to recall that it was mainly concerned with the study of algebraic curves in the complex projective plane, usually approached by projective geometric methods (with or without the use of coordinates). There was a parallel development, with Abel, Jacobi, Weierstrass and Riemann, of the theory of ``algebraic functions'' of one complex variable and their integrals; mathematicians were obviously conscious of the connection between this theory and the geometry of plane algebraic curves and even on occasions were known to ``apply Analysis to Geometry''; but the methods used for the study of algebraic functions were chiefly of a ``transcendental'' nature, even before Riemann (†); this character is still further accentuated in the work of the latter, with the introduction of ``Riemann surfaces'' and arbitrary analytic functions defined on such a surface. Almost immediately after the death of Riemann, Roch and above all Clebsch recognized the possibility of obtaining from the profound results obtained by Riemann's transcendental methods numerous striking applications to the projective geometry of curves, which was of course to incite contemporary geometers to give purely ``geometric'' proofs of these results; this programme, incompletely followed by Clebsch and Gordan, was completed by Brill and M. Noether several years later {[}13{]}, with the aid of the study of systems of variable points on a given curve and auxiliary curves (the ``adjoints'') passing through such systems of points. But even for his contemporaries, Riemann's transcendental methods (and notably his use of topological notions and of ``Dirichlet's principle'') appeared to rest on uncertain foundations; and although Brill and Noether are rather more careful than most contemporary ``synthetic'' geometers (see below p.~593), their geometric-analytic methods are not safe from all reproach. It is essentially to give the theory of plane algebraic curves a solid basis that Dedekind and Weber published in 1882 their great memoir on this subject {[}10 bis{]}: ``The research published below'', they say, ``is intended to lay the foundations of the theory of algebraic functions of one variable, one of Riemann's principal creations, in a way which is at the same time simple, rigorous and entirely general. In earlier research on this subject, in general restrictive hypotheses have been made on the singularities of the functions considered and the would-be exceptional cases are, either mentioned in passing as limiting cases, or entirely neglected. Similarly, certain fundamental theorems on continuity or analyticity are accepted, whose ``evidence'' depends on geometric intuitions of a varied nature'' ({[}10 bis{]}, p.~181) (*).

The essential idea of their work is to model the theory of algebraic functions of one variable on the theory of algebraic numbers as Dedekind had just developed it; to do this, they must first look at it from an ``affine'' point of view (in contrast with their contemporaries, who invariably considered algebraic curves as imbedded in complex projective space): they therefore start with a finite algebraic extension K of the field \(\mathbf{C}(\mathbf{X})\) of rational functions and the ring A of ``integral algebraic functions'' in K, i.e.~the elements of this field which are integral over the polynomial ring C{[}X{]}; their fundamental result, which they obtain without using any topological consideration (*), is that A is a Dedekind domain, to which may be applied mutatis mutandis (and even, as Dedekind and Weber remark, without yet clearly seeing the reason ({[}10{]}, vol.~I, p.~268), in a simpler way) all the results of the ``11th supplement''. Having done this, they prove that their theorems are in fact birationally invariant (in other words, depend only on the field K) and in particular do not depend on the choice of the ``line at infinity'' made at the beginning. What is no doubt still more interesting for us, is that, in order to define the points of the ``Riemann surface'' corresponding to K (and in particular the ``points at infinity'', which cannot correspond to ideals of A), they are led to introduce the notion of place of the field K: They find themselves in the same situation as Gelfand will find himself in 1940 when founding the theory of normed algebras, knowing a set K of elements which are not given to start with as functions and yet one wants to consider as such; and, to obtain the defining set of hypothetical functions, they have for the first time the idea (which Gelfand followed and which has become commonplace through being used at every turn in modern mathematics) of associating with a point x a set E and with a set F of mappings of E to a set G the mapping \(f \mapsto f(x)\) of F to G, in other words of considering, in the expression \(f(x)\) , f as variable and x as fixed, in contrast with the classical tradition. Finally, they have no difficulty, starting from the notion of place, in defining ``positive divisors'' (``Polygon'' in their terminology) which include the ideals of A as particular cases and correspond to the ``systems of points'' of Brill and Noether; but, although they write principal divisors and divisors of differentials as ``quotients'' of positive divisors, they do not give the general definition of divisors and it is only in 1902 that Hensel and Landsberg introduce, by analogy with fractional ideals, this notion which will always embarrass the champions of purely ``geometric'' methods (obliged in spite of themselves to define them with the name ``virtual systems'', but uneasy at not being able to give them a ``concrete'' interpretation).

The same year 1882 also sees appear Kronecker's great memoir awaited for more than 20 years {[}9f{]}. Much more ambitious than the work of Dedekind-

Weber, it is also unfortunately much more vague and obscure. Its central theme is (in modern language) the study of the ideals of a finite integral algebra over one of the polynomial rings \(C[X_{1},\ldots,X_{n}]\) or \(Z[X_{1},\ldots,X_{n}]\) ; Kronecker limits himself a priori to those ideals which are finitely generated (the fact that they all are was only to be proved (for the ideals of \(C[X_{1},\ldots,X_{n}]\) ) some years later by Hilbert in the course of his work on invariants {[}16a{]}). As far as \(C[X_{1},\ldots,X_{n}]\) or \(Z[X_{1},\ldots,X_{n}]\) is concerned, this naturally led to associating with each ideal of one of these rings the ``algebraic variety'' consisting of the zeros common to all the elements of the ideal; and the study of geometry in 2 and 3 dimensions during the 19th century was to lead intuitively to the idea that every variety is a union of a finite number of ``irreducible'' varieties whose ``dimensions'' are not necessarily all the same. It seems that the proof of this fact is the aim Kronecker sets himself, although he nowhere says so explicitly and no definition of ``irreducible variety'' can be found in his memoir, nor of ``dimension''. In fact, he limits himself to indicating summarily how a general elimination method (*) gives, starting with a system of generators of the ideal considered, a finite number of algebraic varieties for each of which, in a suitable coordinate system, a certain number of coordinates are arbitrary and the others are ``algebraic functions'' of them (†). But if it is indeed the decomposition into irreducible varieties at which Kronecker is aiming, it must be recognized that he only arrives there in the elementary case of a principal ideal, where he proves effectively, extending a classical lemma of Gauss on Z{[}X{]} ({[}3{]}, vol.~I, p.~34), that the domains \(C[X_{1},\ldots,X_{n}]\) and \(Z[X_{1},\ldots,X_{n}]\) are factorial; and, in the general case, it is questionable whether Kronecker was in possession of the notion of prime ideal (what he calls ``Primmodulsystem'' is an ideal which is indecomposable as a product of two others ({[}9f{]}, p.~336); this is all the more astonishing as the definition already given by Dedekind in 1871 was perfectly general).

It must however be said that Kronecker's elimination method, suitably applied, certainly leads to the decomposition of an algebraic variety into its irreducible components: this is clearly established by E. Lasker at the beginning of his great memoir of 1905 on polynomial ideals {[}19{]}; he defines correctly the notion of irreducible variety (in \(\mathbf{C}^n\) ) as an algebraic variety V such that a product of two polynomials can only be zero on the whole of the variety V if one of them is and he also gives a definition which is independent of the choice of axes. In the interesting historical considerations he inserts in this work, Lasker shows that he is interested, not only in the purely algebraic tendencies of Kronecker and Dedekind, but also in the problems raised by the geometric methods of the school of Clebsch and M. Noether and notably in the famous theorem proved by the latter in 1873 {[}12{]}. He is essentially concerned, as we would say today, with determining the ideal \(a\) of polynomials of \(\mathbf{C}[\mathbf{X}_1,\ldots ,\mathbf{X}_n]\) which are zero at the points of a given set M of \(\mathbf{C}^n\) ; usually M was the ``algebraic variety'' of zeros common to polynomials \(f_{i}\) finite in number and for a long time it seems that it was accepted (of course without justification) that, at least for \(n = 2\) or \(n = 3\) , the ideal \(a\) was simply generated by the \(f_{i}\) (*). M. Noether had shown that even for \(n = 2\) and for two polynomials \(f_{1},f_{2}\) this is generally false and he had given sufficient conditions for \(a\) to be generated by \(f_{1}\) and \(f_{2}\) . Ten years later, Netto proves that, with no hypothesis on \(f_{1}\) and \(f_{2}\) , a power of \(a\) is always contained in the ideal generated by \(f_{1}\) and \(f_{2}\) {[}15{]}, a theorem which Hilbert generalized in 1893 in his celebrated Nullstellensatz {[}16b{]}. No doubt inspired by this result, Lasker, in his memoir, introduces the general notion of primary ideal ( \(\dagger\) ) in the rings \(\mathbf{C}[\mathbf{X}_1,\dots ,\mathbf{X}_n]\) and \(\mathbf{Z}[\mathbf{X}_1,\dots ,\mathbf{X}_n]\) (after having given for these rings the definition of prime ideal, by transcribing Dedekind's definition) and shows ( \(\ddagger\) ) the existence of a primary decomposition for every ideal in these rings (*). He does not seem to be concerned with questions of uniqueness in this decomposition; it is Macaulay who, a little later {[}21{]} introduces the distinction between ``immersed'' and ``non-immersed'' primary ideals and shows that the latter are determined uniquely, but not the former. It should finally be noted that Lasker also extends his results to the ring of convergent power series in a neighbourhood of a point, by using Weierstrass's ``preparation theorem''. This part of his memoir is no doubt the first place this ring had been considered from a purely algebraic point of view and the methods which Lasker develops on this occasion were strongly to influence Krull when in 1938 he created the general theory of local rings (cf.~{[}29d{]}, p.~204 and passim).

\begin{center}* * *\end{center}

The movement of ideas which will give birth to modern Commutative Algebra begins to take shape about 1910. If the general notion of field was reached by the beginning of the 20th century, in contrast the first work where the general notion of ring is defined is probably that of Fraenkel in 1914 {[}23{]}. At this time, there were already as examples of rings, not only the integral domains of the Theory of Numbers and Algebraic Geometry, but also rings of power series (formal and convergent) and finally algebras (commutative or not) over a base field. However, for the theory of fields as well as that of rings, the catalyst role seems to have been played by Hensel's theory of p-adic numbers, which Fraenkel and also Steinitz {[}20a{]} mention specially as the starting point of their research.

Hensel's first publication on this subject goes back to 1897; he there starts from the analogy shown by Dedekind and Weber between the points of a Riemann surface of an algebraic function field K and the prime ideals of a number field k; he proposes to carry over to the Theory of Numbers ``Puiseux expansions'' (classical from the middle of the 19th century) which, in a neighbourhood of any point of the Riemann surface of K, allow every element \(x \in K\) to be expressed in the form of a convergent series of powers of the ``uniformizer'' at the point considered (a series with only a finite number of terms with negative exponent). Hensel shows similarly that, if p is a prime ideal of k lying above a prime number p, a ``p-adic series'' may be associated with every \(x \in k\) , of the form \(\sum_{i} \alpha_{i} p^{i} \left( \text{or } \sum_{i} \alpha_{i} p^{i/e} \text{ when } p \text{ is ramified over } p \right)\) , the \(\alpha_{i}\) being taken in a given representative system of the field of residues of the ideal p; but his great originality lies in having had the idea of considering such ``expansions'' even when they correspond to no element of k, by analogy with the expansions in integral series of transcendental functions on a Riemann surface {[}18a{]}.

Throughout the rest of his career, Hensel devotes himself to polishing and perfecting little by little his new calculus; and if his manner seems to us hesitant or ponderous, it must not be forgotten that at the beginning at least he had at his disposal none of the topological or algebraic tools of modern mathematics which would have facilitated his task. In his first publications he moreover scarcely speaks of topological notions and on the whole for him the ring of p-adic integers (p a prime ideal in the ring of integers A of a number field k) is, in modern terms, the inverse limit of the rings A/p \(^{n}\) for n increasing indefinitely, in a purely algebraic sense; and to establish the properties of this ring and its field of fractions, it is necessary at each step to use more or less painfully ad hoc arguments (for example to prove that the p-adic integers form an integral domain). The idea of introducing topological notions into a p-adic field does not appear in Hensel's works before 1905 {[}18d{]}; and it is only in 1907, after having published the book where he reexpounds the theory of algebraic numbers according to his ideas {[}18f{]}), that he arrives at the definition and essential properties of p-adic absolute values {[}18e{]}, starting with which he will be able to develop, modeling it on Cauchy's theory, a new ``p-adic analysis'' which he will be able to apply fruitfully in the Theory of Numbers (notably by using the p-adic exponential and logarithm) and whose importance has been growing ever since.

Hensel had well seen, from the beginning, the simplifications his theory brought to classical expositions, by allowing the problems to be ``localized'' and the work to be carried out in a field where not only are the divisibility properties trivial, but also, thanks to the fundamental lemma which he discovered as early as 1902 {[}18c{]}, the study of polynomials whose ``reduced'' polynomials mod p have no multiple roots is reduced to the study of polynomials over a finite field. He had given as early as 1897 {[}18b{]} striking examples of these simplifications, notably on questions related to the discriminant (in particular, a short proof of the criterion he had given a few years earlier for the existence of ``extraordinary divisors''). But for a long time it seems that the p-adic numbers inspired considerable distrust in contemporary mathematicians; a current attitude no doubt towards ideas that are ``too abstract'', but which was also justified in part by the rather excessive enthusiasm of their author (so frequent in mathematics among zealots of new theories). Not content to apply his theory fruitfully to algebraic numbers, Hensel, impressed as all his contemporaries were, by the proofs of the transcendence of e and \(\pi\) and perhaps misled by the adjective ``transcendental'' applied both to numbers and to functions, had come to think that there existed a connection between his p-adic numbers and transcendental real numbers and he had thought for a moment that he had obtained a simple proof of the transcendence of e and even of \(e^{e}\) ({[}18d{]}, p.~556) (*).

Soon after 1910, the situation changes, with the rising of the next generation, influenced by the ideas of Fréchet and F. Riesz on topology and by those of Steinitz on algebra, and from the start devoted to ``abstraction''; it will know how to assimilate and put in their true place Hensel's works. As early as 1913, Kürschak {[}22{]} gives a general definition of the notion of absolute value, recognizes the importance of ultrametric absolute values (of which the p-adic absolute value was an example), proves (by modelling the proof on the case of real numbers) the existence of the completion of a field with respect to an absolute value and above all shows generally the possibility of extending an absolute value to any algebraic extension of the given field. But he had not seen that the ultrametric character of an absolute value was already revealed in the prime field; this point was established by Ostrowski, to whom also is due the determination of all the absolute values on the field Q and the fundamental theorem characterizing fields with a non-ultrametric absolute value as subfields of C {[}24{]}. In the years from 1920 to 1935, the theory will be completed by a more detailed study of absolute values which are not necessarily discrete, including amongst others the examination of various circumstances which arise in passing to an algebraic or transcendental extension (Ostrowski, Deuring, F. K. Schmidt); on the other hand, in 1931, Krull introduces and studies the general notion of valuation {[}29b{]} which will be greatly used in the years that follow by Zariski and his school of Algebraic Geometry (†). We must also mention here, although it lies outside our scope, the deeper studies on the structure of complete valued fields and complete local rings, which date from the same period (Hasse-Schmidt, Witt, Teichmüller, I. Cohen).

\begin{center}* * *\end{center}

The work of Fraenkel mentioned above (p.~594) only treated a very special

type of ring (Artinian with only a single prime ideal, which is moreover assumed to be principal). With the exception of Steinitz's work on fields {[}20a{]}, the first important works on the study of general commutative rings are E. Noether's two great memoirs on ideal theory: that of 1921 {[}25a{]}, devoted to primary decomposition, which takes up again in all generality and completes on many points the results of Lasker and Macaulay; and that of 1927 characterizing Dedekind domains axiomatically {[}25b{]}. Just as Steinitz had shown for fields, it is seen in these memoirs how a small number of abstract ideas, such as the notion of irreducible ideal, the chain conditions and the idea of an integrally closed domain (the last two, as we have seen, already brought to light by Dedekind) can by themselves lead to general results which seemed inextricably bound up with results of pure computation in the cases where they had previously been known.

With these memoirs of E. Noether, joined to the slightly later works of Artin-van der Waerden on divisorial ideals {[}31{]} and Krull relating these ideals to essential valuations {[}29b{]}, the long study of the decomposition of ideals started a century earlier (*) is thus complete, at the same time as modern Commutative Algebra is being inaugurated.

The innumerable later research works on Commutative Algebra are grouped most easily according to several important directions of development:

\textbf{(A) Local rings and topologies}\par

Although the germ was contained in all the earlier works on the Theory of Numbers and Algebraic Geometry, the general idea of localization came to light very slowly. The general notion of ring of fractions is only defined in 1926 by H. Grell, a pupil of E. Noether, and only for integral domains {[}28{]}; its extension to more general rings will only be given in 1944 by C. Chevalley for Noetherian rings and in 1948 by Uzkov in the general case. Until about 1940, Krull and his school are practically alone in using in general arguments the consideration of the local rings \(\mathbf{A}_{\mathfrak{p}}\) of an integral domain \(\mathbf{A}\) ; these rings will only begin to appear explicitly in Algebraic Geometry with the works of Chevalley and Zariski starting in 1940 (†).

The general study of local rings themselves only begins in 1938 with Krull's great memoir {[}29d{]}. The most important results of this work concern dimension theory and regular rings, of which we shall not speak here; but here for the first time appears the completion of any arbitrary Noetherian local ring and also a still imperfect form of the graded ring associated with a local ring (*); the latter will only be defined about 1948 by P. Samuel {[}36{]} and independently in research on Algebraic Topology by Leray and H. Cartan. Krull, in the above mentioned work, hardly uses topological language; but already in 1928 {[}29a{]}, he had proved that, in a Noetherian ring A, the intersection of the powers of an ideal a is the set of \(x \in A\) such that \(x(1 - a) = 0\) for some \(a \in a\) ; it is easily deduced from this that, for every ideal m of A, the m-adic topology on A induces on an ideal a the m-adic topology on a; in his memoir of 1938, Krull completes this result by proving that in a Noetherian local ring every ideal is closed. These theorems were soon afterwards extended by Chevalley to Noetherian semi-local rings and then by Zariski to the rings which bear his name {[}33b{]}; to Chevalley also goes back the introduction of ``linear compactness'' in topological rings, as also the determination of the structure of complete semi-local rings {[}32b{]}.

\textbf{(B) Passage from the local to the global}\par

Since Weierstrass, an analytic function of one variable (and in particular an algebraic function) has habitually been associated with the set of its ``expansions'' at all the points of the Riemann surface where it is defined. In the introduction to his book on the Theory of Numbers ({[}18f{]}, p.~V), Hensel similarly associates with each element of an algebraic number field k the set of elements corresponding to it in the completions of k with respect to all the absolute values on \(k (\dagger)\) . It may be said that it is this point of view which, in modern Commutative Algebra, has replaced the decomposition formula of an ideal as a product of prime ideals (extending in a certain sense Kummer's initial point of view). Hensel's remark amounts implicitly to embedding k in the product of all its completions; this is what Chevalley does explicitly in 1936 with his theory of ``idèles'' {[}32a{]}, which perfects earlier analogous ideas of Prüfer and von Neumann (the latter confining themselves to embedding \(k\) in the product of its p-adic completions) (*). Although this is somewhat outside our scope, it is important to mention here that, thanks to an appropriate topology on the group of idèles, all the techniques of locally compact groups (including Haar measure) can thus be very effectively applied to the Theory of Numbers.

In a more general context, Krull's theorem {[}29b{]} characterizing an integrally closed domain as an intersection of valuation rings (which amounts also to embedding the domain under consideration in a product of valuation rings) often facilitates the study of these rings, although the method is only really tractable for essential valuations of Krull domains. Moreover Krull frequently exhibits {[}29e{]} (quite elementary) examples of the ``passage from the local to the global'' method consisting of showing a property of an integral domain A by verifying it for the ``localized'' rings \(A_{p}\) of A at all its prime ideals (†); more recently, Serre perceived that this method works for arbitrary commutative rings A, that it is applicable also to A-modules and that it is even sufficient often to ``localize'' at the maximal ideals of A (Chapter II, § 3, Theorem 1): a point of view which is closely connected with ideas about ``spectra'' and sheaves defined over these spectra (see below, p.~602).

\textbf{(C) Integers and integral closure}\par

We have seen that the notion of algebraic integer, first introduced for number fields, had already been extended by Kronecker and Dedekind to algebraic function fields, although in this case it might appear rather artificial (not corresponding to a projective notion). E. Noether's memoir of 1927, followed by the work of Krull starting in 1931, showed the interest that these notions present for more general rings ( \(\ddagger\) ). Krull in particular is responsible for the theorems on the lifting of prime ideals to integral algebras {[}29c{]}, as also for extending the theory of decomposition and inertia groups of Dedekind-Hilbert {[}29b{]}. As for E. Noether, we owe to her the general formulation of the normalization lemma (*) (from which follows amongst other things Hilbert's Nullstellensatz) as also the first general criterion (transcribing the classical arguments of Kronecker and Dedekind) for the integral closure of an integral domain to be finite over that domain.

Finally, it should be pointed out here that one of the reasons for the present importance of the notion of integrally closed domain is due to Zariski's studies on algebraic varieties; he discovered that ``normal'' varieties (that is those whose local rings are integrally closed domains) are distinguished by particularly pleasant properties, notably the fact that they have no ``singularity of codimension 1''; and it has then been seen that analogous phenomena are true for ``analytic spaces''. Therefore ``normalization'' (that is the operation which, for the local rings of a variety, consists of taking their integral closures) has become a powerful weapon in the arsenal of modern Algebraic Geometry.

\textbf{(D) The study of modules and the influence of Homological Algebra}\par

One of the striking characteristics of the work of E. Noether and W. Krull in Algebra is the tendency to ``linearization'', extending the analogous development given to field theory by Dedekind and Steinitz; in other words, ideals are considered above all as modules and so all the constructions of Linear Algebra (quotient, product and more recently tensor product and formation of homomorphism modules) are brought to bear on them, producing in general modules which are no longer ideals. It is thus quickly seen that in many questions (for commutative or non-commutative rings), interest should not be confined to the study of ideals of a ring A, but on the contrary the theorems should be stated in general for A-modules (sometimes subjected to certain finiteness conditions).

The intervention of Homological Algebra has strongly reinforced the above tendency, since this branch of Algebra is essentially concerned with questions of a linear nature. We shall not retrace its history here; but it is interesting to point out that several fundamental notions of Homological Algebra (such as that of projective module and that of Tor functor) came into being on the occasion of a close study of the behaviour of modules over a Dedekind domain relative to the tensor product, a study undertaken by H. Cartan in 1948.

Conversely, it could be foreseen that the new classes of modules introduced naturally by Homological Algebra as ``universal annihilators'' of the Ext functors (projective modules and injective modules) and the Tor functors (flat modules) would throw new light on Commutative Algebra. It happens that chiefly projective modules and still more flat modules have shown themselves useful: the importance of the latter arises above all from the remark, made first by Serre {[}38b{]}, that localization and completion introduce flat modules naturally, thus ``explaining'' in a much more satisfactory way the properties of these operations already known and rendering them much easier to use. It should moreover be mentioned (as we shall see in later chapters) that the applications of Homological Algebra are far from being limited to this and that it is playing a more and more important role in Algebraic Geometry.

\textbf{(E) The notion of spectrum}\par

The most recent in date of the new notions of Commutative Algebra has a complex history. Hilbert's spectral theorem introduced ordered sets of orthogonal projectors of a Hilbert space, forming a ``Boolean algebra'' (or rather a Boolean lattice) (*), in one-to-one correspondence with a Boolean lattice of classes of measurable subsets (for a suitable measure) of \(\mathbf{R}\) . No doubt his earlier work on operators on Hilbert spaces, about 1935, led M. H. Stone to study Boolean lattices generally and notably to look for ``representations'' of them by subsets of a set (or classes of subsets with respect to a certain equivalence relation). He observes that a Boolean lattice becomes a commutative ring (moreover of a very special type), if multiplication is defined on it by \(xy = \inf(x, y)\) and addition by \(x + y = \sup(\inf(x, y'), \inf(x', y))\) . In the particular case where the Boolean lattice in question is the set \(\mathfrak{P}(\mathbf{X})\) of all subsets of a finite set \(\mathbf{X}\) , it is immediately seen that the elements of \(\mathbf{X}\) are in a natural one-to-one correspondence with the maximal ideals of the corresponding ``Boolean'' ring; and Stone obtains precisely his general representation theorem for a Boolean lattice by similarly considering the set of maximal ideals of the corresponding ring and associating with each element of the Boolean lattice the set of maximal ideals which contains it {[}30a{]}.

On the other hand, the set of both open and closed subsets of a topological space was a well-known classical example of a Boolean lattice. In a second paper {[}30b{]}, Stone showed that in fact every Boolean lattice is also isomorphic to a Boolean lattice of this nature. For this it was of course necessary to define a topology on the set of maximal ideals of a ``Boolean'' ring; which was very simply accomplished by taking as closed sets for each ideal a the set of maximal ideals containing a.

We shall not speak here of the influence of these ideas on Functional Analysis, where they played an important role in the birth of the theory of normed algebras developed by I. Gelfand and his school. But in 1945, Jacobson observes {[}34{]} that the process of defining a topology, invented by Stone, can in fact be applied to any ring A (commutative or not) provided the set of ideals taken is not the set of maximal ideals but the set of two-sided ``primitive'' ideals (i.e.~the two-sided ideals b such that A/b is a primitive ring); for a commutative ring, these of course turn out to be the maximal ideals. On his part, Zariski, in 1944 {[}33a{]}, uses an analogous method to define a topology on the set of places of an algebraic function field. However, these topologies remained for the majority of algebraists mere curiosities, by reason of the fact that they are not usually Hausdorff and a quite understandable repugnance was felt about working on such unusual objects. This distrust was only overcome when A. Weil showed, in 1952, that every algebraic variety can be given a natural topology of the above type and that this topology allows the definition, in perfect analogy with the case of differentiable or analytic manifolds, of the notion of fibre bundle {[}37{]}; soon afterwards, Serre had the idea of extending to these varieties thus topologized the theory of coherent sheaves, thanks to which the topology renders in the case of ``abstract'' varieties the same services as the usual topology when the base field is C, notably as far as applying the methods of Algebraic Topology is concerned ({[}38a{]} and {[}38b{]}).

From then on it was natural to use this geometric language throughout Commutative Algebra. It was quickly seen that considering maximal ideals is usually insufficient to obtain useful assertions (*) and that the adequate notion is that of the set of prime ideals of the ring topologized in the same manner. With the introduction of the notion of spectrum, there now exists a dictionary allowing every theorem of Commutative Algebra to be expressed in a geometric language very close to that of the Algebraic Geometry of the Weil-Zariski period; which has moreover immediately brought about a considerable enlargement of the scope of the latter, so that Commutative Algebra is scarcely more than the most elementary part of it {[}39{]}.
